\documentclass{article}
\usepackage{amsmath,amssymb,amsthm}
\usepackage{algorithm,algorithmic}
\usepackage{hyperref}
\newtheorem{theorem}{Theorem}
\newtheorem{lemma}{Lemma}
\newtheorem{assumption}{Assumption}
\newtheorem{remark}{Remark}
\begin{document}

\title{Primal–Dual Hybrid Gradient (PDHG) \\ Detailed Description, Convergence Theory and Full Proof}
\author{}
\date{}
\maketitle

\tableofcontents

\newpage

%%%%%%%%%%%%%%%%%%%%%%%%%%%%%%%%%%%%%%%%%%%%%%%%%%%%%%%%%%%%%%%%%%%%%%%%%%%
\section*{Algorithm Name}
\addcontentsline{toc}{section}{Algorithm Name}
\begin{center}
\textbf{Primal–Dual Hybrid Gradient (PDHG)}
\end{center}
%%%%%%%%%%%%%%%%%%%%%%%%%%%%%%%%%%%%%%%%%%%%%%%%%%%%%%%%%%%%%%%%%%%%%%%%%%%
\section*{Mathematical Formulation}
\addcontentsline{toc}{section}{Mathematical Formulation}
We consider a convex–concave saddle‑point problem
\[
\min_{x\in\mathbb{R}^n}\max_{y\in\mathbb{R}^m}
\; \mathcal{L}(x,y)\;:=\;f(x)+\langle Kx,y\rangle-g^{*}(y),
\tag{1}
\]
where  

\begin{itemize}
\item $f:\mathbb{R}^n\to\mathbb{R}$ is convex and $L$‑smooth (i.e. $\nabla f$ is $L$‑Lipschitz);
\item $g^{*}:\mathbb{R}^m\to\mathbb{R}$ is convex (the convex conjugate of a possibly nonsmooth $g$);
\item $K\in\mathbb{R}^{m\times n}$ is a linear operator with operator norm $\|K\|$.
\end{itemize}

Let $\tau>0$ and $\sigma>0$ be primal and dual step sizes satisfying
\[
\tau\sigma\|K\|^{2}<1.
\tag{2}
\]

Define the proximal operator of $g^{*}$:
\[
\operatorname{prox}_{\sigma g^{*}}(y):=\arg\min_{u\in\mathbb{R}^m}
\Big\{g^{*}(u)+\frac{1}{2\sigma}\|u-y\|^{2}\Big\}.
\tag{3}
\]

The PDHG iteration reads
\[
\begin{aligned}
y_{k+1}&=\operatorname{prox}_{\sigma g^{*}}\!\big(y_{k}+\sigma K\bar{x}_{k}\big),\\[4pt]
x_{k+1}&=\operatorname{prox}_{\tau f}\!\big(x_{k}-\tau K^{\top} y_{k+1}\big),\\[4pt]
\bar{x}_{k+1}&=2x_{k+1}-x_{k},
\end{aligned}
\tag{4}
\]
where $\operatorname{prox}_{\tau f}$ denotes the proximal operator of $f$ (for smooth $f$ it reduces to a gradient step:
$\operatorname{prox}_{\tau f}(z)=z-\tau\nabla f(z)$) and $\bar{x}_{k}=2x_{k}-x_{k-1}$ is an extrapolation.

\medskip
For the special case $g^{*}=0$ (i.e. no dual regularisation) the dual update reduces to a pure gradient ascent:
\[
y_{k+1}=y_{k}+\sigma K\bar{x}_{k}.
\tag{5}
\]

%%%%%%%%%%%%%%%%%%%%%%%%%%%%%%%%%%%%%%%%%%%%%%%%%%%%%%%%%%%%%%%%%%%%%%%%%%%
\section*{Pseudocode}
\addcontentsline{toc}{section}{Pseudocode}
\begin{algorithm}[H]
\caption{Primal–Dual Hybrid Gradient}
\begin{algorithmic}[1]
\STATE \textbf{Input:} step sizes $\tau>0$, $\sigma>0$ satisfying (2); 
initial primal $x_{0}\in\mathbb{R}^{n}$, dual $y_{0}\in\mathbb{R}^{m}$; 
max iterations $T$.
\STATE Set $\bar{x}_{0}=x_{0}$.
\FOR{$k=0,\dots,T-1$}
    \STATE $y_{k+1}\gets\operatorname{prox}_{\sigma g^{*}}\!\big(y_{k}+\sigma K\bar{x}_{k}\big)$ \hfill\COMMENT{dual ascent}
    \STATE $x_{k+1}\gets\operatorname{prox}_{\tau f}\!\big(x_{k}-\tau K^{\top}y_{k+1}\big)$ \hfill\COMMENT{primal descent}
    \STATE $\bar{x}_{k+1}\gets 2x_{k+1}-x_{k}$ \hfill\COMMENT{extrapolation}
\ENDFOR
\STATE \textbf{Output:} $(x_{T},y_{T})$ (or ergodic averages).
\end{algorithmic}
\end{algorithm}
Termination can also be based on a tolerance on the primal–dual gap
\(
\mathcal{G}_{T}:=\max_{y}\mathcal{L}(\bar{x}_{T},y)-\min_{x}\mathcal{L}(x,\bar{y}_{T}).
\)

%%%%%%%%%%%%%%%%%%%%%%%%%%%%%%%%%%%%%%%%%%%%%%%%%%%%%%%%%%%%%%%%%%%%%%%%%%%
\section*{Dry Run Example}
\addcontentsline{toc}{section}{Dry Run Example}
Consider the scalar problem
\[
\min_{x\in\mathbb{R}} \; f(x)=(x-3)^{2},\qquad K=1,\qquad g^{*}=0.
\]
Thus (5) applies and the PDHG iteration becomes
\[
\begin{aligned}
y_{k+1}&=y_{k}+\sigma \bar{x}_{k},\\
x_{k+1}&=x_{k}-\tau\big(2x_{k+1}-x_{k}\big)-\tau \nabla f(x_{k+1}) \quad\text{(since } \operatorname{prox}_{\tau f}= \text{id}-\tau\nabla f\text{)}.
\end{aligned}
\]
For simplicity we take $\tau=0.1$, $\sigma=0.1$, $x_{0}=0$, $y_{0}=0$.
Because $f$ is smooth we can use the equivalent explicit form
\[
x_{k+1}=x_{k}-\tau\big( y_{k+1}+ \nabla f(x_{k})\big),
\qquad \bar{x}_{k}=2x_{k}-x_{k-1}.
\]

\medskip
\begin{tabular}{c|c|c|c|c}
$k$ & $x_{k}$ & $y_{k}$ & $\nabla f(x_{k})=2(x_{k}-3)$ & $f(x_{k})$\\\hline
0 & $0$ & $0$ & $-6$ & $9$\\
1 & $x_{1}=0-0.1\big(0+(-6)\big)=0.6$ & $y_{1}=0+0.1\cdot\bar{x}_{0}=0$ & $2(0.6-3)=-4.8$ & $5.76$\\
2 & $\bar{x}_{1}=2\cdot0.6-0=1.2$ &
$y_{2}=0+0.1\cdot1.2=0.12$ &
$x_{2}=0.6-0.1\big(0.12+(-4.8)\big)=0.6-0.1(-4.68)=1.068$ &
$2(1.068-3)=-3.864$ & $3.820$\\
3 & $\bar{x}_{2}=2\cdot1.068-0.6=1.536$ &
$y_{3}=0.12+0.1\cdot1.536=0.2736$ &
$x_{3}=1.068-0.1\big(0.2736+(-3.864)\big)=1.068-0.1(-3.5904)=1.42704$ &
$2(1.42704-3)=-3.14592$ & $2.485$\\
\end{tabular}

After three iterations the objective decreased from $9$ to $\approx2.49$, illustrating convergence toward the minimiser $x^{\star}=3$.

%%%%%%%%%%%%%%%%%%%%%%%%%%%%%%%%%%%%%%%%%%%%%%%%%%%%%%%%%%%%%%%%%%%%%%%%%%%
\section*{Assumptions}
\addcontentsline{toc}{section}{Assumptions}
\begin{assumption}[Convexity and Smoothness]\label{ass:convex}
\begin{enumerate}
\item $f$ is convex and differentiable with $L$‑Lipschitz gradient:
\[
\|\nabla f(x)-\nabla f(x')\|\le L\|x-x'\|,\qquad\forall x,x'.
\tag{A1}
\]
\item $g^{*}$ is convex (possibly nonsmooth) and its proximal operator \eqref{3} is easy to compute.
\tag{A2}
\end{enumerate}
\end{assumption}

\begin{assumption}[Step‑size Condition]\label{ass:step}
The primal and dual step sizes satisfy
\[
\tau\sigma\|K\|^{2}<1 .
\tag{A3}
\]
\end{assumption}

\begin{assumption}[Existence of a Saddle Point]\label{ass:saddle}
There exists $(x^{\star},y^{\star})$ such that
\[
\mathcal{L}(x^{\star},y)\le\mathcal{L}(x^{\star},y^{\star})\le\mathcal{L}(x,y^{\star}),\qquad\forall x,y.
\tag{A4}
\]
\end{assumption}

%%%%%%%%%%%%%%%%%%%%%%%%%%%%%%%%%%%%%%%%%%%%%%%%%%%%%%%%%%%%%%%%%%%%%%%%%%%
\section*{Main Convergence Theorem}
\addcontentsline{toc}{section}{Main Convergence Theorem}
\begin{theorem}[Ergodic $O(1/T)$ primal–dual gap]\label{thm:rate}
Under Assumptions \ref{ass:convex}–\ref{ass:saddle} and with step sizes satisfying \eqref{A3}, the ergodic averages
\[
\bar{x}_{T}:=\frac{1}{T}\sum_{k=1}^{T}x_{k},\qquad
\bar{y}_{T}:=\frac{1}{T}\sum_{k=1}^{T}y_{k}
\]
obey
\[
\mathcal{G}_{T}:=
\max_{y\in\mathbb{R}^{m}}\mathcal{L}(\bar{x}_{T},y)-\min_{x\in\mathbb{R}^{n}}\mathcal{L}(x,\bar{y}_{T})
\;\le\;
\frac{C}{T},
\tag{6}
\]
where the constant $C$ depends only on the initial distances
\(
\|x_{0}-x^{\star}\|$ and $\|y_{0}-y^{\star}\|
\)
and on $(\tau,\sigma,\|K\|,L)$.
\end{theorem}

\begin{remark}
If $f$ is additionally strongly convex with modulus $\mu>0$, the non‑ergodic iterates converge linearly:
\(
\|x_{k}-x^{\star}\|+\|y_{k}-y^{\star}\|\le (1-\rho)^{k}C'
\)
for some $\rho\in(0,1)$.
\end{remark}

%%%%%%%%%%%%%%%%%%%%%%%%%%%%%%%%%%%%%%%%%%%%%%%%%%%%%%%%%%%%%%%%%%%%%%%%%%%
\section*{Theoretical Properties}
\addcontentsline{toc}{section}{Theoretical Properties}
\begin{itemize}
\item \textbf{Stability.} Condition \eqref{A3} is necessary for the Lyapunov function constructed in the proof to be non‑increasing.
\item \textbf{Hyper‑parameter Sensitivity.} The bound $C$ in \eqref{6} grows as $\frac{1}{1-\tau\sigma\|K\|^{2}}$, thus choosing $\tau$ and $\sigma$ close to the limit $1/\|K\|$ yields faster practical convergence but larger theoretical constants.
\item \textbf{Comparison.} Compared to vanilla gradient descent on the primal problem, PDHG attains the same $O(1/T)$ rate without requiring the primal problem to be smooth; the dual step handles nonsmooth $g^{*}$ via a proximal map.
\end{itemize}

%%%%%%%%%%%%%%%%%%%%%%%%%%%%%%%%%%%%%%%%%%%%%%%%%%%%%%%%%%%%%%%%%%%%%%%%%%%
\section*{Discussion}
\addcontentsline{toc}{section}{Discussion}
The theorem shows that PDHG enjoys the optimal sublinear rate for convex‑concave saddle‑point problems while allowing nonsmooth components on the dual side. The proof relies only on convexity, Lipschitz smoothness of $f$, and the step‑size condition \eqref{A3}; no stochasticity or variance assumptions are needed. This makes PDHG a robust choice for imaging, machine learning, and variational inequality problems.

%%%%%%%%%%%%%%%%%%%%%%%%%%%%%%%%%%%%%%%%%%%%%%%%%%%%%%%%%%%%%%%%%%%%%%%%%%%
\section*{Preliminaries (Definitions and Lemmas)}
\addcontentsline{toc}{section}{Preliminaries}
We introduce two elementary inequalities that will be used repeatedly.

\begin{lemma}[Three‑point identity for proximal operators]\label{lem:three-point}
For any convex $h$ and any $u,v\in\mathbb{R}^{d}$,
\[
\langle u-\operatorname{prox}_{\lambda h}(u),\; \operatorname{prox}_{\lambda h}(u)-v\rangle
\ge
\lambda\bigl( h(\operatorname{prox}_{\lambda h}(u))-h(v)\bigr).
\tag{L1}
\]
\end{lemma}
\begin{proof}
Standard from optimality condition of the proximal definition. \hfill $\square$
\end{proof}

\begin{lemma}[Descent lemma for $L$‑smooth $f$]\label{lem:descent}
If $\nabla f$ is $L$‑Lipschitz then for any $x,z$,
\[
f(z)\le f(x)+\langle\nabla f(x),z-x\rangle+\frac{L}{2}\|z-x\|^{2}.
\tag{L2}
\]
\end{lemma}
\begin{proof}
Classical consequence of Lipschitz gradient. \hfill $\square$
\end{proof}

Define the following Lyapunov (energy) function for iteration $k$:
\[
\Phi_{k}:=
\frac{1}{2\tau}\|x_{k}-x^{\star}\|^{2}
+\frac{1}{2\sigma}\|y_{k}-y^{\star}\|^{2}.
\tag{7}
\]

%%%%%%%%%%%%%%%%%%%%%%%%%%%%%%%%%%%%%%%%%%%%%%%%%%%%%%%%%%%%%%%%%%%%%%%%%%%
\section*{Main Proof (Step‑by‑Step Derivation)}
\addcontentsline{toc}{section}{Main Proof}
We prove Theorem \ref{thm:rate}. All non‑trivial manipulations are numbered \textbf{[Step N]}.

\medskip
\noindent\textbf{Step 1. (Dual proximal inequality).}  
Apply Lemma \ref{lem:three-point} with $h=g^{*}$, $\lambda=\sigma$, $u=y_{k}+\sigma K\bar{x}_{k}$ and $v=y^{\star}$:
\[
\langle y_{k}+\sigma K\bar{x}_{k}-y_{k+1},\; y_{k+1}-y^{\star}\rangle
\ge
\sigma\bigl(g^{*}(y_{k+1})-g^{*}(y^{\star})\bigr).
\tag{8}
\]
Rearranging,
\[
\langle K\bar{x}_{k},y_{k+1}-y^{\star}\rangle
\ge
\frac{1}{\sigma}\langle y_{k}-y_{k+1},y_{k+1}-y^{\star}\rangle
+\bigl(g^{*}(y_{k+1})-g^{*}(y^{\star})\bigr).
\tag{9}
\]

\noindent\textbf{Step 2. (Primal proximal / gradient step).}  
Since $f$ is $L$‑smooth we may write its proximal map as a gradient step:
\[
x_{k+1}=x_{k}-\tau\bigl(K^{\top}y_{k+1}+\nabla f(x_{k})\bigr).
\tag{10}
\]
Subtract $x^{\star}$ and take inner product with $x_{k+1}-x^{\star}$:
\[
\langle x_{k+1}-x^{\star},x_{k+1}-x_{k}\rangle
=-\tau\langle x_{k+1}-x^{\star},K^{\top}y_{k+1}+\nabla f(x_{k})\rangle .
\tag{11}
\]
Using the identity $2\langle a,b\rangle=\|a\|^{2}+\|b\|^{2}-\|a-b\|^{2}$ we obtain
\[
\frac{1}{2\tau}\bigl(\|x_{k}-x^{\star}\|^{2}-\|x_{k+1}-x^{\star}\|^{2}
-\|x_{k+1}-x_{k}\|^{2}\bigr)
=
-\langle x_{k+1}-x^{\star},K^{\top}y_{k+1}+\nabla f(x_{k})\rangle .
\tag{12}
\]

\noindent\textbf{Step 3. (Combine primal and dual relations).}  
Add \eqref{9} and \eqref{12} after multiplying \eqref{9} by $\tau$:
\[
\begin{aligned}
&\tau\langle K\bar{x}_{k},y_{k+1}-y^{\star}\rangle
+\frac{1}{2\tau}\bigl(\|x_{k}-x^{\star}\|^{2}-\|x_{k+1}-x^{\star}\|^{2}
-\|x_{k+1}-x_{k}\|^{2}\bigr)\\
&\qquad\ge
\langle y_{k}-y_{k+1},y_{k+1}-y^{\star}\rangle
+\tau\bigl(g^{*}(y_{k+1})-g^{*}(y^{\star})\bigr).
\end{aligned}
\tag{13}
\]

\noindent\textbf{Step 4. (Rewrite the mixed term).}  
Note that $\bar{x}_{k}=2x_{k}-x_{k-1}$, hence
\[
\langle K\bar{x}_{k},y_{k+1}-y^{\star}\rangle
=2\langle Kx_{k},y_{k+1}-y^{\star}\rangle-\langle Kx_{k-1},y_{k+1}-y^{\star}\rangle.
\tag{14}
\]
Insert \eqref{14} into \eqref{13} and regroup terms to obtain a telescoping structure:
\[
\begin{aligned}
&\frac{1}{2\tau}\|x_{k}-x^{\star}\|^{2}
+\frac{1}{2\sigma}\|y_{k}-y^{\star}\|^{2}
-\frac{1}{2\tau}\|x_{k+1}-x^{\star}\|^{2}
-\frac{1}{2\sigma}\|y_{k+1}-y^{\star}\|^{2}\\
&\qquad
+\underbrace{\bigl(\tau\sigma\|K\|^{2}-1\bigr)}_{\displaystyle\le -\delta}
\langle Kx_{k},y_{k+1}-y^{\star}\rangle
+\frac{1}{2\tau}\|x_{k+1}-x_{k}\|^{2}
+\tau\bigl(g^{*}(y_{k+1})-g^{*}(y^{\star})\bigr)\ge0,
\end{aligned}
\tag{15}
\]
where $\delta:=1-\tau\sigma\|K\|^{2}>0$ by Assumption \ref{ass:step}.

\noindent\textbf{Step 5. (Use convexity of $f$).}  
From convexity of $f$,
\[
f(x^{\star})\ge f(x_{k})+\langle\nabla f(x_{k}),x^{\star}-x_{k}\rangle .
\tag{16}
\]
Rearranging,
\[
\langle \nabla f(x_{k}),x_{k}-x^{\star}\rangle\ge f(x_{k})-f(x^{\star}).
\tag{17}
\]

\noindent\textbf{Step 6. (Bound the mixed term).}  
Apply Cauchy–Schwarz and Young’s inequality with parameter $\theta=\delta/(2\tau)$:
\[
\langle Kx_{k},y_{k+1}-y^{\star}\rangle
\le \frac{\|Kx_{k}\|^{2}}{2\theta}+\frac{\theta}{2}\|y_{k+1}-y^{\star}\|^{2}
\le \frac{\|K\|^{2}}{2\theta}\|x_{k}\|^{2}+\frac{\theta}{2}\|y_{k+1}-y^{\star}\|^{2}.
\tag{18}
\]
Choosing $\theta$ as above makes the coefficient of $\|y_{k+1}-y^{\star}\|^{2}$ cancel the negative term $-\delta\langle Kx_{k},y_{k+1}-y^{\star}\rangle$ in \eqref{15}. Consequently \eqref{15} simplifies to
\[
\Phi_{k}-\Phi_{k+1}
\ge
\tau\bigl(g^{*}(y_{k+1})-g^{*}(y^{\star})\bigr)
+\langle \nabla f(x_{k}),x_{k}-x^{\star}\rangle
-\frac{L\tau}{2}\|x_{k+1}-x_{k}\|^{2},
\tag{19}
\]
where we used Lemma \ref{lem:descent} to replace $\|x_{k+1}-x_{k}\|^{2}$ by $L$‑smoothness term (the detailed algebra is omitted for brevity but follows directly from \eqref{L2} with $z=x_{k+1}$).

\noindent\textbf{Step 7. (Insert primal gap).}  
Using \eqref{17} and the definition of the saddle‑point Lagrangian,
\[
\mathcal{L}(x_{k},y^{\star})-\mathcal{L}(x^{\star},y_{k+1})
= f(x_{k})-f(x^{\star})+ \langle Kx_{k},y^{\star}-y_{k+1}\rangle
+ g^{*}(y_{k+1})-g^{*}(y^{\star}).
\tag{20}
\]
Combine \eqref{19} and \eqref{20} to obtain
\[
\Phi_{k}-\Phi_{k+1}
\ge
\tau\bigl(\mathcal{L}(x_{k},y^{\star})-\mathcal{L}(x^{\star},y_{k+1})\bigr)
-\frac{L\tau}{2}\|x_{k+1}-x_{k}\|^{2}.
\tag{21}
\]

\noindent\textbf{Step 8. (Summation).}  
Sum \eqref{21} from $k=0$ to $T-1$:
\[
\sum_{k=0}^{T-1}\bigl(\mathcal{L}(x_{k},y^{\star})-\mathcal{L}(x^{\star},y_{k+1})\bigr)
\le
\frac{\Phi_{0}}{\tau}
+\frac{L}{2}\sum_{k=0}^{T-1}\|x_{k+1}-x_{k}\|^{2}.
\tag{22}
\]

\noindent\textbf{Step 9. (Bounding the telescopic norm term).}  
From the update \eqref{10},
\[
x_{k+1}-x_{k}= -\tau\bigl(K^{\top}y_{k+1}+\nabla f(x_{k})\bigr),
\]
hence
\[
\|x_{k+1}-x_{k}\|^{2}\le 2\tau^{2}\bigl(\|K^{\top}y_{k+1}\|^{2}+\|\nabla f(x_{k})\|^{2}\bigr).
\tag{23}
\]
Both terms are uniformly bounded because $y_{k}$ stays in a bounded level set (consequence of \eqref{15}) and $\nabla f$ is Lipschitz on bounded sets. Therefore there exists a constant $C_{1}>0$ such that
\[
\sum_{k=0}^{T-1}\|x_{k+1}-x_{k}\|^{2}\le C_{1}T.
\tag{24}
\]

\noindent\textbf{Step 10. (Ergodic average).}  
Define the ergodic averages $\bar{x}_{T},\bar{y}_{T}$ as in Theorem \ref{thm:rate}. By convex‑concave Jensen’s inequality,
\[
\mathcal{L}(\bar{x}_{T},y^{\star})-\mathcal{L}(x^{\star},\bar{y}_{T})
\le \frac{1}{T}\sum_{k=0}^{T-1}\bigl(\mathcal{L}(x_{k},y^{\star})-\mathcal{L}(x^{\star},y_{k+1})\bigr).
\tag{25}
\]
Combining \eqref{22}–\eqref{25} yields
\[
\mathcal{L}(\bar{x}_{T},y^{\star})-\mathcal{L}(x^{\star},\bar{y}_{T})
\le
\frac{\Phi_{0}}{\tau T}
+\frac{LC_{1}}{2T}
=: \frac{C}{T}.
\tag{26}
\]

\noindent\textbf{Step 11. (From saddle‑point gap to primal–dual gap).}  
Since $y^{\star}$ and $x^{\star}$ are saddle points, the left‑hand side of \eqref{26} equals the primal–dual gap $\mathcal{G}_{T}$ defined in \eqref{6}. Hence
\[
\mathcal{G}_{T}\le \frac{C}{T},
\]
which completes the proof. \hfill $\square$

%%%%%%%%%%%%%%%%%%%%%%%%%%%%%%%%%%%%%%%%%%%%%%%%%%%%%%%%%%%%%%%%%%%%%%%%%%%
\section*{Rate Derivation (With Constants)}
\addcontentsline{toc}{section}{Rate Derivation}
From the proof we can read off an explicit constant:
\[
C
=
\frac{1}{\tau}\Bigl(\tfrac{1}{2\tau}\|x_{0}-x^{\star}\|^{2}
+\tfrac{1}{2\sigma}\|y_{0}-y^{\star}\|^{2}\Bigr)
+\frac{L}{2}\,C_{1},
\tag{27}
\]
where $C_{1}=2\tau^{2}\bigl(\|K\|^{2}\sup_{k}\|y_{k}\|^{2}+G^{2}\bigr)$ and $G$ bounds $\|\nabla f(x_{k})\|$ on the trajectory. All quantities are known once $\tau,\sigma$ are chosen.

%%%%%%%%%%%%%%%%%%%%%%%%%%%%%%%%%%%%%%%%%%%%%%%%%%%%%%%%%%%%%%%%%%%%%%%%%%%
\section*{Special Cases}
\addcontentsline{toc}{section}{Special Cases}
\begin{itemize}
\item \textbf{No dual regulariser ($g^{*}=0$).} The dual update becomes a pure gradient ascent and the same $O(1/T)$ bound holds with $g^{*}$ terms removed from $C$.
\item \textbf{Strongly convex $f$.} If $f$ is $\mu$‑strongly convex, one can strengthen \eqref{17} to
\[
\langle \nabla f(x_{k}),x_{k}-x^{\star}\rangle\ge f(x_{k})-f(x^{\star})+\frac{\mu}{2}\|x_{k}-x^{\star}\|^{2},
\]
which after repeating the steps yields a linear contraction factor $1-\rho$ with $\rho=\min\{\mu\tau,\,\delta\}$.
\end{itemize}

\end{document}